% Propietario conservado: /Users/ruben/Documents/ChatGPT/jueces y controles/REVISION_ARGUMENTAL_APERTURAS_CIERRES_20260915/II/manuscript_es/sections/10b_boltzmann.tex
\section{Constante de Boltzmann, información y transducción térmica}
\label{ii:sec:ii-boltzmann}

\subsection{Estructura dimensional de la transducción}

La definición estructural de la constante de Boltzmann utiliza las
rectas de temperatura y de energía. La transducción positiva
\[
 k_B:\mathcal L_\Theta\longrightarrow\mathcal L_E
\]
asigna a una temperatura su escala energética. En la normalización
trítica, la sección de temperatura del ciclo satisface
\[
 k_B\Theta_{\rm clk}=\frac{E_P}{\log3}.
\]
El miembro derecho queda fijado por la acción y el período previamente
construidos. El factor \(\log3\) procede del contenido informacional
de una distribución uniforme sobre los tres estados locales.

Sean \(u_E,u_\Theta\) bases positivas y escribamos
\(E_P=\epsilon_Pu_E\), \(\Theta_{\rm clk}=\theta_{\rm clk}u_\Theta\)
y \(k_B(u_\Theta)=b\,u_E\). Entonces
\[
 b\,\theta_{\rm clk}=\frac{\epsilon_P}{\log3}.
\]
Para \(u'_E=s_Eu_E\) y \(u'_\Theta=s_\Theta u_\Theta\),
las coordenadas se transforman mediante
\[
 b'=\frac{s_\Theta}{s_E}b,\qquad
 \theta'_{\rm clk}=\frac{\theta_{\rm clk}}{s_\Theta},\qquad
 \epsilon'_P=\frac{\epsilon_P}{s_E}.
\]
La identidad energética permanece covariante. Esta formulación reúne
la estructura individual del transductor y la normalización de la
sección sobre la que actúa.

\begin{remark}[Coordenada térmica y elección de bases]
El transductor y la identidad energética anterior son covariantes.
Un valor numérico individual de \(k_B\) está referido a las bases positivas
\(u_E,u_\Theta\); su transformación es \(b'=(s_\Theta/s_E)b\).
Por tanto, la selección interna de una coordenada numérica requiere,
además, una construcción que fije esas bases. La identidad energética
no selecciona por sí sola tal elección. La coordenada metrológica
no define la acción ni el período.
\end{remark}

La distinción ayuda a interpretar la entropía. El funcional
\(s(\rho)=-\operatorname{Tr}(\rho\log\rho)\) es adimensional.
La entropía dimensional es \(S(\rho)=k_Bs(\rho)\), entendida en
la recta correspondiente. La cantidad \(\Theta S\) pertenece a
la recta de energía. Así se conserva el sentido matemático de las
igualdades que reúnen información, temperatura y acción.

En particular, la información de continuación
\eqref{ii:eq:ii-entropia-continuaciones} admite la publicación dimensional
\(S_n^{\rm gen}(x)=k_Bs_n(x)\). Las ramas admisibles y su medida
determinan \(s_n\); el transductor fija la escala de entropía. Para
una distribución uniforme sobre tres alternativas, el valor es
\(k_B\log3\). Cuando las prolongaciones poseen pesos diferentes,
la expresión ponderada conserva esa diferencia estructural. La
constante de Boltzmann puede definirse aquí por su función: es el
coeficiente dimensional que transporta la información de la
distribución admitida a la recta de entropía y vincula, bajo una
temperatura declarada, esa información con energía.

La acción interviene antes de la publicación térmica. La sección
\(\hbar_a\) convierte la frecuencia de la evolución en energía;
la sección \(h_a\) expresa la acción acumulada durante una vuelta.
La igualdad \(k_B\Theta_{{\rm clk},a}\log3=h_a/T_{108}\)
reúne ambas funciones: energía del período e información local de
la dinámica trítica. Una identificación termodinámica conserva
además el estado estadístico, la temperatura y la operación física
de intercambio o borrado que se estudia.
